\startsection[title={Biorreatores}]
  Biorreatores fornecem ambientes favoráveis para o crescimento de culturas de microrganismos e podem ser otimizados para a produção de determinado produto bioquímico, desde que haja acesso a mecanismos de controle das variáveis físicas envolvidas, o que geralmente é oferecido por reatores projetados especificamente para essa finalidade, contrariamente aos bioreatores de ocorrencia natural, onde a manipulaçao consistente de variaveis nao acontece com facilidade ou é completamente mitigada - a exemplo, uma poça de agua a ceu aberto com atividade microbiana em que se deseja controlar a temperatura. A necessidade de um projeto de biorreator decorre da responsabilidade de garantir a reprodutibilidade dos processos bioquímicos desenvolvidos ao longo de uma cadeia produtiva. A possibilidade de controlar algumas variáveis \emdash ou até mesmo todas aquelas avaliadas como impactantes para a conversão final \emdash é uma característica indispensável para aplicações industriais de rotas produtivas desenvolvidas em laboratório. Sendo que uma vez determinadas as condições de operação na etapa experimental, basta ajustar-se as variáveis extensivas do processo para uma escala industrial.

  Segundo \cite[authoryears][ref::erickson_0], existem diversas formas de classificar biorreatores, seja pelo tipo de alimentação ou pelos produtos finais obtidos. Existem três tipos principais de alimentação: batelada, alimentação alternada e alimentação contínua. A alimentação em batelada ocorre quando o substrato é introduzido no sistema apenas no início do processo, sem manipulações externas posteriores na quantidade de fonte de carbono. A alimentação alternada consiste em uma estratégia de distribuição da fonte de carbono seguindo um cronograma, diferenciando-se da primeira pela possibilidade de introdução de substrato durante o experimento. Por fim, a alimentação contínua caracteriza-se pela introdução constante de alimento ao longo de todo o processo, representando uma tentativa de reproduzir um sistema estacionário a partir da manutenção da concentração de substrato no biorreator. A \in{Figura}[fig:simple_bioreactor] mostra um examplo de volume de controle que pode ser pensado como um bioreator.

  \startplacefigure[title={Ilustração de biorreator genérico}, reference={fig:simple_bioreactor}]
    \externalfigure[image/simple_bioreactor.png][height=200pt]
  \stopplacefigure

  O bioreator representado na \in{Figura}[fig:simple_bioreactor] tem uma corrente de entrada \math{F_{i}} e de saida \math{F_{o}}, com suas devida composições. Nota-se o conjunto de variáveis (coordenadas) \math{\left\lparen T, pH, p\chemical{O_{2}}\right\rparen}, temperatura, pH, e pressão parcial de oxigênio respectivamente, sugerindo uma monitoração dessas variáveis, porque elas são os parametros de interesse para o processo bioquímico apresentado. Portanto, um bioreator pode ser considerado um tanque, ou volume de controle projetado especificamente para adaptar condições otimas para uma cultura de microorganismos metabolisar o alimento presente segundo uma estequimetria ideal, ou proxima disso \cite[authoryear][ref::erickson_0].

  Quando tratamos de processos aeróbicos, é necessário fornecer oxigênio para o desenvolvimento da cultura. Nessa perspectiva, boa parte dos reatores destinados ao cultivo de organismos aeróbicos necessita de agitação para uniformizar a concentração de oxigênio em todo o meio, o que é naturalmente dificultado pela baixa solubilidade do oxigênio em água. Isso pode ser realizado por meio de reatores de mistura, colunas de ar dissolvido, membranas ou pelo aumento da área de interface líquido-gás, o que não ocorre necessariamente em processos anaeróbicos ou em sistemas envolvendo culturas mistas \cite[authoryear][ref::erickson_0]. Essa última opção tem como proposta estimular a competição entre diferentes variantes presentes com base na seleção natural. Portanto, os tipos de cultura influenciam diretamente no projeto físico do biorreator.

  Outra condição de operação que influencia drasticamente tanto o projeto físico quanto o modelo matemático de um reator é a fase dos reagentes e produtos. Em fase líquida, a cinética bioquímica depende da concentração dos reagentes e, em alguns casos, também da concentração dos produtos. Para muitas reações bioquímicas, essa abordagem cinética pode ser tratada, com boa aproximação, considerando dependência de primeira ordem em baixas concentrações e de ordem zero em concentrações elevadas \cite[authoryear][ref::erickson_0]. Em fase sólida, a cinética bioquímica torna-se mais dependente da disposição espacial do substrato sólido e do contato superficial deste substrato com o microrganismo. Dessa forma, uma solubilidade substancial dos reagentes é essencial para garantir adequada difusão no meio e, consequentemente, a formação do produto de interesse \cite[authoryear][ref::bhargav_0]. Em alguns sistemas, utilizam-se bioenzimas para hidrolisar o substrato em reagentes mais solúveis, aumentando a conversão por unidade de tempo.

  O controle das variáveis físicas envolvidas é importante para a predição da resposta do sistema em relação a qualquer perturbação de entrada, bem como para a manipulação do sistema, de forma que os resultados de réplicas experimentais sejam os mais semelhantes possíveis dentro das limitações de sistemas não determinísticos. Entre as variáveis que podem ser medidas em um biorreator estão temperatura, pressão, pH, potência de agitação, fluxo de gases e líquidos, concentração de gases, concentração de biomassa e tensão interfacial \cite[authoryear][ref::erickson_0].
\stopsection

